\documentclass[12pt]{article}
\usepackage[T1]{fontenc}
\usepackage[utf8]{inputenc}
\usepackage{lmodern}
\usepackage{textcomp}
\usepackage{enumitem}
\usepackage{geometry}
\geometry{margin=1in}
\begin{document}
\begin{center}
Answer Key - Religious Studies - K-12 Level Assessment 3 \
\small (Answers are ordered exactly as in the question paper.)
\end{center}

\section*{Multiple Choice}
\begin{enumerate}[label=\arabic*.]
\item \textbf{Answer:} D
\\ \textit{Options:} A) Septuagint; B) Talmud; C) Yeshiva; D) Gemarah
\\ \textit{Note:} The correct answer is Gemarah because it refers to the body of Jewish rabbinical analysis and commentary on the Mishnah, which together with the Mishnah forms the Talmud.
\item \textbf{Answer:} D
\\ \textit{Options:} A) Haggadah; B) Halakhah; C) Brit; D) Bris
\\ \textit{Note:} The correct answer is "Bris" because in Yiddish, the term "Bris" specifically refers to the covenant of circumcision, which is a significant religious ritual in Judaism symbolizing the covenant between God and the Jewish people.
\item \textbf{Answer:} A
\\ \textit{Options:} A) Disciple; B) Faithful; C) Devotee; D) Enlightened
\\ \textit{Note:} The Punjabi word "Sikh" translates to "disciple," reflecting the core principle of the Sikh faith, which emphasizes following the teachings of the Guru and seeking spiritual guidance.
\item \textbf{Answer:} B
\\ \textit{Options:} A) Parvati; B) Ganesha; C) Vishnu; D) Sita
\\ \textit{Note:} Ganesha is considered the most beloved elephant-headed Hindu deity because he symbolizes wisdom, prosperity, and the removal of obstacles, making him a widely revered figure in Hindu worship and culture.
\item \textbf{Answer:} B
\\ \textit{Options:} A) Presbyteries; B) Sees; C) Churches; D) Synods
\\ \textit{Note:} The term "sees" refers to the five principal centers of early Christianity, known as the Pentarchy, which included the sees of Rome, Constantinople, Alexandria, Antioch, and Jerusalem, highlighting their significance in the governance and doctrinal authority of the Church.
\end{enumerate}

\section*{True / False}
\begin{enumerate}[label=\arabic*.]
\setcounter{enumi}{5}
\item \textbf{Answer:} True
\\ \textit{Options:} A) True; B) False
\\ \textit{Note:} Pure Land Buddhism emphasizes that individuals can attain salvation and rebirth in the Pure Land through their faith in Amida Buddha's compassion and the recitation of his name, reflecting the belief in his saving grace.
\item \textbf{Answer:} True
\\ \textit{Options:} A) True; B) False
\\ \textit{Note:} The statement is true because the title 'Mahatma', which translates to 'Great Soul' in Sanskrit, symbolizes Gandhi's profound spiritual influence and ethical guidance in the Indian independence movement.
\item \textbf{Answer:} False
\\ \textit{Options:} A) True; B) False
\\ \textit{Note:} Corinth was not the primary location for mystery cults in ancient Greece; instead, places like Eleusis and Delphi were more renowned for their mystery religions, attracting significant numbers of worshippers seeking spiritual enlightenment.
\item \textbf{Answer:} False
\\ \textit{Options:} A) True; B) False
\\ \textit{Note:} The Tripitaka is actually the traditional scripture of Buddhism, consisting of three "baskets" of teachings, while the term "Three gems" refers to the core principles of Buddhism: the Buddha, the Dharma, and the Sangha.
\item \textbf{Answer:} False
\\ \textit{Options:} A) True; B) False
\\ \textit{Note:} Mudras are not clothing but rather symbolic hand gestures used in Buddhist practices to convey spiritual meanings and enhance meditation.
\end{enumerate}

\section*{Long Form Response}
\begin{enumerate}[label=\arabic*.]
\setcounter{enumi}{10}
\item \textbf{Answer:} The belief that kingship in Sumer descends from heaven is significant as it establishes a divine legitimacy for rulers, suggesting that their authority is ordained by the gods. This concept influenced Sumerian society by creating a strong connection between religion and governance, where kings were viewed as intermediaries between the gods and the people. As a result, the king's decisions and laws were often justified by divine will, reinforcing the social hierarchy and encouraging loyalty among citizens. This belief also shaped rituals and practices, as rulers engaged in religious ceremonies to maintain the favor of the gods, further intertwining governance and spirituality in Sumerian life.
\\ \textit{Note:} correct because it highlights how the divine origin of kingship in Sumer legitimized rulers' authority, intertwined religion with governance, and reinforced social hierarchies, thereby fostering loyalty among the populace.
\item \textbf{Answer:} Thich Nhat Hanh, a prominent Buddhist monk, emphasizes the concept of self-sacrifice as a compassionate act aimed at alleviating the suffering of others. In Buddhism, self- sacrifice is not about losing oneself but rather about understanding the interconnectedness of all beings and acting selflessly to promote peace and happiness. This theme influences contemporary Buddhist practices by encouraging mindfulness, community service, and the cultivation of loving-kindness, inspiring individuals to prioritize the welfare of others. In today's society, these values manifest in various forms, such as volunteer work and social activism, reflecting a commitment to building a more compassionate world. Ultimately, Thich Nhat Hanh's teachings remind us that self- sacrifice can lead to greater harmony and understanding within our communities.
\\ \textit{Note:} correct because it articulates Thich Nhat Hanh's interpretation of self-sacrifice as an expression of compassion rooted in the interconnectedness of all beings, and it further connects this understanding to contemporary Buddhist practices that prioritize mindfulness, community service, and loving-kindness.
\end{enumerate}

\vfill
\noindent\textit{End of Answer Key}
\end{document}